% Generated draft: requires mathematical and source review.
\documentclass[11pt]{article}
\usepackage[margin=25mm]{geometry}
\usepackage{fontspec}
\setmainfont{Latin Modern Roman}
\usepackage{amsmath,amssymb,mathtools,unicode-math}
\setmathfont{Latin Modern Math}
\usepackage{enumitem,needspace,xurl,hyperref}
\setlength{\parindent}{0pt}
\setlength{\parskip}{5pt}
\setlength{\emergencystretch}{3em}
\newcommand{\N}{\mathbb{N}}
\newcommand{\Z}{\mathbb{Z}}
\newcommand{\Q}{\mathbb{Q}}
\newcommand{\R}{\mathbb{R}}
\newcommand{\C}{\mathbb{C}}
\newcommand{\F}{\mathbb{F}}
\newcommand{\A}{\mathbb{A}}
\newcommand{\PP}{\mathbb{P}}
\newcommand{\E}{\mathbb{E}}
\newcommand{\eps}{\varepsilon}
\newcommand{\dd}{\,\mathrm{d}}
\begin{document}
\section*{UnsolvedMath ID: 30003849}
\textbf{Source problem code:} catalogue entry
\subsection{Definitions and mathematical statement}
The supplied record combines three independent questions and is marked \texttt{unrecoverable}; it does not provide a verified unified formulation. The questions below are therefore retained separately, with unresolved conventions stated explicitly.

Write \(V_{\mathrm{br}}\) and \(F_{\mathrm{br}}\) for the braided Thompson groups named in the source, \(nV\) for the named Brin--Thompson groups, and \(T\) for Thompson's group appearing in the Stein--Farley question. The supplied record gives no presentations or constructions identifying these groups or the particular cube complex precisely.

A group \(G\) is \emph{Hopfian} if every surjective endomorphism of \(G\) is injective. For a coefficient ring \(R\), let \(H_k(G;R)\) denote group homology with trivial coefficients. The usual integral meaning of \emph{acyclic} is
\[
H_k(G;\Z)=0\qquad\text{for every integer }k\geq 1.
\]
The source does not specify homology coefficients for its acyclicity questions; the integral interpretation below is consequently provisional.

\begin{enumerate}
\item Determine separately whether \(V_{\mathrm{br}}\) and \(F_{\mathrm{br}}\) are Hopfian. Precisely, for each named group \(G\), determine whether
\[
\forall\varphi\in\operatorname{End}(G),\qquad
\varphi(G)=G\ \Longrightarrow\ \ker\varphi=\{1\}.
\]
\item Determine whether \(V_{\mathrm{br}}\) is acyclic, and determine the corresponding acyclicity question for the Brin--Thompson groups \(nV\). Under the conventional integral interpretation, these ask whether
\[
H_k(V_{\mathrm{br}};\Z)=0\qquad(k\geq 1)
\]
and whether
\[
H_k(nV;\Z)=0\qquad(k\geq 1)
\]
for the intended values of \(n\). The source does not specify the range of \(n\), or whether it seeks a universal assertion or a classification by \(n\); no such quantifier is imposed here.
\item Can the Stein--Farley cube complex for \(T\) be used to recover the computation attributed to Ghys--Sergiescu of
\[
H_*(T;\Q)?
\]
This is a question about obtaining a specified computation by a specified method. The supplied record neither states that computation nor identifies the cube complex, its action, or the required sense of recovery, so it does not determine a more precise theorem to prove.
\end{enumerate}

The supplied verification note reports that the official report contains three independent numbered questions merged by this extraction. Their individual numbers and full statements are absent here. These formulations require human mathematical review and do not constitute a recovered clean statement.
\subsection{Short English statement}
The record combines separate questions about Hopficity of \(V_{\mathrm{br}}\) and \(F_{\mathrm{br}}\), acyclicity of \(V_{\mathrm{br}}\) and the Brin--Thompson groups \(nV\), and recovering the Ghys--Sergiescu rational homology computation for \(T\) using a Stein--Farley cube complex. The acyclicity coefficients, intended range of \(n\), and precise computational target remain unresolved. The supplied open label is dated curation, not an independently verified conclusion.
\subsection{Sources}
\begin{itemize}
\item Supplied local registry record: \texttt{data/unsolvedmath/problems.json}, permanent UnsolvedMath ID \(30003849\), source code \texttt{catalogue entry}, title ``Homology and Hopficity of Braided Thompson Groups,'' in the Oberwolfach Reports Open Problems collection. The record gives proposed year 2018 and no named proposer.
\item Supplied report citation: \emph{Cohomological and Metric Properties of Groups of Homeomorphisms of R} (2019), Oberwolfach Reports, pp.\ 1579--1633, DOI \texttt{10.4171/owr/2018/26}. \url{https://doi.org/10.4171/owr/2018/26}. The supplied report citation is authoritative provenance, but its contents have not been independently consulted here. The supplied verification note states that the extraction merges three independent numbered questions and marks the statement \texttt{unrecoverable}.
\item Supplied literature source: Bux--Flossner--Meyer--Witzel, ``braided Thompson groups.'' \url{https://arxiv.org/abs/1210.2931}.
\item Supplied literature source: Fournier-Facio--Lodha--Zaremsky, ``braided Thompson groups.'' \url{https://doi.org/10.2140/agt.2024.24.1601}.
\item Supplied literature source: Oberwolfach Report 26/2018. \url{https://ems.press/content/serial-article-files/46748?nt=1}.
\item Literature assessment checked at \texttt{2026-08-21}: the supplied curation labels the record open and states that known type \(F_\infty\) and bounded-cohomology results for related braided groups do not settle the Hopfian, acyclicity, Brin--Thompson, or Stein--Farley questions in this record. This is supplied dated curation derived from the OWR material, not an independently verified literature assessment; its status and cited claims require independent verification.
\item The rational homology computation is attributed in the supplied statement to Ghys--Sergiescu. No separate bibliographic reference or formula for that computation is supplied.
\end{itemize}
\end{document}
